\chapter{\rm\bfseries Experimental Setup and Methodology}
\label{ch:methods}

\section{Hardware Configuration}

\subsection{Robot Manipulator}

A Franka Emika FR3 7-DOF cobot served as the manipulator throughout this work.
The arm was operated through the Franka Control Interface (FCI) using
\texttt{libfranka 0.18.1} and the ROS\,2 Humble control stack. Cartesian
pose tracking was performed by a custom \texttt{fr3\_pose\_controller}
maintained by the McGill Applied Dynamics Lab at
\url{https://github.com/McGill-Applied-Dynamics-Lab/adl-ros2},
parameterised in \path{config/controllers/fr3_pose/default.yaml}. The controller accepts a
sequence of (pose, twist) waypoints with associated time stamps and
interpolates between them using a quintic polynomial; for the in-place
re-orientation and the final precision descent phases, a cosine half-cycle
S-curve interpolator (the so-called ``sine descent'' implemented in
\path{run_trial.py::_safe_traj}) is used to ensure zero start- and end-state
velocities.

\subsection{Gripper Variants Tested}

\begin{table}[ht]
\centering
\caption[Gripper variants]{Six gripper variants benchmarked. All finray
variants share the same 3D-printed inclined-seat mount; only the seat
angle changes between the $+2.5^{\circ}$ baseline and the $+5.0^{\circ}$
variant. The parallel-jaw variants do not use an inclined seat.}
\label{tab:variants}
\begin{tabular}{llll}
\toprule
Variant & Type & Key parameter & Seat angle \\
\midrule
\texttt{parallel\_jaw\_stock}    & Rigid two-finger (Franka Hand)      & ---                    & --- \\
\texttt{parallel\_jaw\_TPU}      & Rigid two-finger + TPU 90A pad      & 2.0\,mm pad thickness  & --- \\
\texttt{finray\_18}              & Finray-effect 3D-printed            & 18\,mm finger thickness & $+2.5^{\circ}$ \\
\texttt{finray\_18\_model2}      & Second print of \texttt{finray\_18} & 18\,mm finger thickness & $+2.5^{\circ}$ \\
\texttt{finray\_26}              & Finray-effect 3D-printed            & 26\,mm finger thickness & $+2.5^{\circ}$ \\
\texttt{finray\_26\_5deg}        & Finray-26 on 5$^{\circ}$-seat mount & 26\,mm finger thickness & $+5.0^{\circ}$ \\
\bottomrule
\end{tabular}
\end{table}

The finray fingers and the parallel-jaw compliant pad were 3D-printed from
Bambu Lab TPU 90A (Shore-A hardness 90) on a Bambu Lab X1 Carbon printer
using the vendor-default TPU 90A profile in Bambu Studio (0.4\,mm nozzle,
default layer height, default infill, PEI-textured build plate).

\paragraph{Baseline inclined-seat mount.}
All finray variants are mounted on the Franka Hand via an inclined seat
that introduces a fixed $+2.5^{\circ}$ tilt about the gripper $y$ axis
(the gripping direction). The seat is a 3D-printed adapter and is
present in \emph{every} variant tested in this thesis unless otherwise
noted. Variants whose name contains \texttt{\_5deg} substitute a
$+5.0^{\circ}$ seat in place of the default $+2.5^{\circ}$ seat; all
other variants (\texttt{finray\_18}, \texttt{finray\_18\_model2},
\texttt{finray\_26}, and the parallel-jaw variants) inherit the
$+2.5^{\circ}$ baseline. The $+2.5^{\circ}$ baseline seat angle was chosen so that the two
opposing finray fingers converge inward toward each other rather than
splaying apart, closing the residual slit that would otherwise remain
between the finger tips when the gripper is fully closed. This ensures
firm contact between the fingers along their full length during a
successful grasp. Figure~\TODO{ref} shows the mounting stack
(flange $\to$ seat $\to$ finger) and defines the seat-tilt convention.

\paragraph{Consequence for reported ``$0^{\circ}$'' cells.}
Cells reported in Chapter~\ref{ch:results} as $0^{\circ}$ tilt refer
to \emph{zero commanded runtime tilt}; the physical gripper is still
inclined at the baseline seat angle relative to the flange. All
runtime tilt values (e.g. $\pm 2.5^{\circ}$ about $x$ or $y$) are
applied \emph{on top of} the baseline seat tilt.

\subsection{Cloth Specimen}
\label{sec:cloth}

\TODO{insert photograph of the cloth specimen and short caption
describing material, thickness, weave, and dimensions.}

\subsection{Test Surface}
\label{sec:test-surface}

\TODO{insert photograph of the test bench in situ and short caption
describing the table material, flatness, and cloth-edge orientation
relative to the gripper $y$ axis.}

\section{Reference Frame and Tool-Centre Calibration}

\subsection{Frame Convention}
\label{sec:frames}

Three frames are relevant: the world (robot base) frame $\mathcal{W}$,
the gripper flange frame $\mathcal{F}$, and the gripper tip frame $\mathcal{T}$.
The Franka URDF defines $\mathcal{F}$ at the mechanical hand mounting flange;
the Franka API \texttt{robot.end\_effector\_pose} reports the pose of
$\mathcal{F}$ in $\mathcal{W}$. The tip frame $\mathcal{T}$ is located along
the local $+z$ axis of $\mathcal{F}$ by a per-variant offset $z_{\text{TCP}}$,
which depends on the physical gripper (finger length, mounting fixtures, pads).

For a gripper at orientation $R \in SO(3)$, the world-frame relationship is
\begin{equation}
  \mathbf{p}_\mathcal{T} = \mathbf{p}_\mathcal{F} + R \begin{bmatrix} 0 \\ 0 \\ z_{\text{TCP}} \end{bmatrix}.
  \label{eq:tip-from-flange}
\end{equation}
For the rigid-down-pointing nominal orientation $R_0 = R_x(180^{\circ})$,
the tip lies directly below the flange and
$\mathbf{p}_\mathcal{T} = \mathbf{p}_\mathcal{F} - z_{\text{TCP}} \hat{\mathbf{z}}$.

The gripper-frame convention adopted here is: the gripping axis (along which
the fingers close) is $\mathcal{F}$'s local $y$ axis; the gripper is
mirror-symmetric about the $xz$ plane; consequently, rotation of the gripper
about the local $y$ axis is rotation about the geometrically asymmetric
axis, while rotation about $x$ is rotation about the symmetric axis.

\subsection{TCP Offset Calibration}
\label{sec:tcp-calib}

The $z_{\text{TCP}}$ value differs by variant because different fingers
attach to the same flange. To calibrate, each variant was driven to a
common probe position $(x_0, y_0)$ on a known rigid table and a force-feedback
touch-off was performed: the gripper descended slowly in $z$ until the
external force in $z$ exceeded a $-3$\,N threshold (debounced for three
consecutive 50\,Hz samples). The flange $z$ at the latch moment was
recorded as $z_{\text{tcp}}^{\text{measured}}$.

Taking the rigid \texttt{parallel\_jaw\_stock} variant as the URDF reference
($z_{\text{TCP}} \equiv 0$ by definition), the per-variant TCP offset is
\begin{equation}
  z_{\text{TCP}}^{(v)} = z_{\text{tcp}}^{(v),\text{measured}}
                       - z_{\text{tcp}}^{\text{stock},\text{measured}}.
  \label{eq:tcp-calib}
\end{equation}

Table~\ref{tab:tcp-calib} summarises the calibrated values used throughout
this work.

\begin{table}[ht]
\centering
\caption[Calibrated TCP offsets]{Calibrated per-variant TCP offsets from
differential touch-off against the stock parallel jaw.}
\label{tab:tcp-calib}
\begin{tabular}{lrr}
\toprule
Variant & $z_{\text{tcp}}^{\text{measured}}$ (mm) & $z_{\text{TCP}}^{(v)}$ (mm) \\
\midrule
\texttt{parallel\_jaw\_stock}    & 37.871 & 0.000 \\
\texttt{parallel\_jaw\_TPU}      & 39.916 & 2.044 \\
\texttt{finray\_18}              & 131.842 & 93.971 \\
\texttt{finray\_26}              & 132.782 & 94.911 \\
\texttt{finray\_26\_5deg}        & 130.094 & 92.223 \\
\bottomrule
\end{tabular}
\end{table}

The reproducibility of the touch-off was measured as the across-tap standard
deviation $\sigma_z$ in mm. For \texttt{finray\_26\_5deg}, $\sigma_z =
183$\,\textmu m owing to a wobble component introduced by the inclined
seat; the other four variants had $\sigma_z \leq 102$\,\textmu m.

\section{Force Measurement}

External forces $\mathbf{F}_{\text{ext}}$ at the flange were estimated by the
Franka internal model from joint torques, available via the \texttt{libfranka}
RobotState API at 1\,kHz and subsampled to 30\,Hz for session logging.
A free-air bias capture was performed at the start of every session: the
gripper was held at a known pose 30--80\,mm above the cloth surface and
30 samples were averaged over 1 second to provide $F_z^{\text{bias}}$.
This bias was recorded in each session JSON but \emph{not} subtracted from
the raw force time series at acquisition time; downstream analyses subtract
it where appropriate.

\section{Tilt-Aware Pose Targeting}
\label{sec:tilt-pose}

When the gripper is tilted, the flange position required to place the tip at
a fixed probe location $(x_0, y_0, z_{\text{surf}}^{\text{tip}})$ shifts
laterally. From equation~\eqref{eq:tip-from-flange}, the required flange
position is
\begin{equation}
  \mathbf{p}_\mathcal{F}^{*} = \mathbf{p}_\mathcal{T}^{\text{target}} -
                              R\,\begin{bmatrix} 0 \\ 0 \\ z_{\text{TCP}} \end{bmatrix}.
  \label{eq:flange-target-tilted}
\end{equation}

This relation is implemented in the trial controller \path{run_trial.py}:
the target orientation $R = R_{\text{tilt}} \circ R_0$ is constructed at
session start by composing the requested tilt rotation
$R_{\text{tilt}} = R_y(\theta_y) R_x(\theta_x)$ (extrinsic, world frame)
onto the base orientation $R_0 = R_x(180^{\circ})$, and the flange position
target is updated accordingly. The orientation is then held constant
through all subsequent descent and grasp phases, so the tip XY remains at
$(x_0, y_0)$ throughout the trial.

For a $+2.5^{\circ}$ tilt about $x$ on the \texttt{finray\_26} variant
($z_{\text{TCP}} = 94.911$\,mm), the predicted flange offset is
\begin{equation}
  \Delta y_\mathcal{F} = -z_{\text{TCP}}\sin(2.5^{\circ})
                       \approx -4.14\,\mathrm{mm}.
\end{equation}
Chapter~\ref{ch:experiments} reports the empirical lateral offsets measured under
load and the residuals from this geometric prediction, which we attribute to
finger compliance under press.

\section{Data Pipeline}

\subsection{Per-Trial Recording}

Each trial directory \path{trial_NNN/} contains:
\begin{itemize}
  \item \path{config.yaml} --- snapshot of the variant config at trial time;
  \item \path{run_config.json} --- depth, tilt, ground-truth source,
        speeds, gripper mode;
  \item \path{pose_wrench.csv} --- high-rate (30\,Hz) actual + commanded
        pose, RPY, wrench, gripper width per row, tagged with the current
        trial stage.
\end{itemize}

\subsection{Per-Session Recording}

Each session directory \path{depth_*/} contains in addition:
\begin{itemize}
  \item \path{depth_summary.json} --- header metadata + array of per-trial
        verdicts (\texttt{success}, \texttt{fail\_reason}); operator notes;
  \item \path{ground_truth_used.json} --- snapshot of the GT JSON loaded;
  \item \path{session_log_*.npz} --- numpy archive of the 30\,Hz background
        sample stream, with stage labels and tap indices.
\end{itemize}

\subsection{Verdict Recording and Outcome Classification}
\label{sec:verdict-recording}

Trial outcomes were classified into three categories rather than a binary
success/failure. After each trial the operator was prompted to distinguish
\emph{success} (the gripper held the cloth throughout the lift stage) from
\emph{no-grasp} (the trial executed to completion and the stationary
force sample was valid, but the cloth slipped from the gripper during
or after the lift). Trials in which the FCI control daemon aborted the
motion mid-trial were recorded as \emph{failed trials} and excluded
from all aggregate analyses; the two most common triggers were:
\begin{itemize}
  \item the Franka safety reflex firing on excessive predicted joint
        force (\texttt{auto\_failed\_excess\_force}), which typically
        occurred when the parallel-jaw variants were commanded to
        depths beyond their narrow force-limited working range;
  \item the Franka safety reflex firing on acceleration discontinuity
        between two consecutive Cartesian trajectory segments, which
        was observed intermittently during the chunked touch-off
        descent and mitigated through the trajectory-emission-floor
        modifications documented in Section~\ref{sec:tcp-calib}.
\end{itemize}
Both trigger classes leave the trial's \texttt{success} field as
\texttt{null} in the per-trial JSON; no-grasp trials, by contrast,
carry \texttt{success = false} with a valid force record and are
retained for engagement-force analysis.

\section{Trial Procedure}

A single trial proceeds:
\begin{enumerate}
  \item Open gripper (Franka Hand or finray release width);
  \item Bulk descent: smooth multi-waypoint trajectory from the lifted pose
        down to a fine clearance ($+5$\,mm above surface), filtered to
        exclude any upward leg (preventing direction reversals at zero-twist
        interior waypoints);
  \item Settle: 2\,s pause at the fine-clearance pose;
  \item Final approach: cosine-S-curve descent from the fine clearance to
        the press target $z_{\text{surf}} - d$, where $d$ is the trial press
        depth;
  \item Grasp: force-controlled close (Franka grasp action at 80\,N or
        20\,N nominal, with $\pm 100$\,mm $\epsilon$ to disable the
        position-acceptance check that would otherwise cause auto-release
        on a sub-spec final width);
  \item Settle: 1\,s pause;
  \item Stationary sample: 30 samples of $\mathbf{F}_{\text{ext}}$
        captured over 1\,s during the \texttt{grasp\_stationary\_sample}
        stage;
  \item Lift to $\sim +80$\,mm above surface for visual inspection;
  \item Operator verdict prompt;
  \item Open gripper and reset for next trial.
\end{enumerate}

\section{Statistical Treatment}

For each (variant, depth, tilt) cell, the contact force is reported as
\begin{equation}
  \bar{F}_z = \frac{1}{N}\sum_{i=1}^{N} F_z^{(i)}, \qquad
  \sigma_{F_z} = \sqrt{\frac{1}{N}\sum (F_z^{(i)} - \bar{F}_z)^2},
\end{equation}
where $N$ is the total number of stationary samples pooled across all
clean trials in the cell. Trials with null verdict (FCI crash) are
excluded. Trial-to-trial variance is the dominant source of dispersion;
the within-trial sample variance ($\sim 0.1$\,N) is small compared to it.
